
\subsection{Theory Directions}
In our research were able to highlight several directions that one could take.
We mainly highlight to directions: \textbf{HF-PPAD} in  \textbf{PPAD}
and topological representations of directed graphs

\subsubsection{\textbf{HF-PPAD} in  \textbf{PPAD}}

In our sections \ref{sec:pc-plus-eol} and \ref{sec:sc-excess}, made references to
how Kleene logic can be incorporated with \textbf{PPAD} problems. For
the \textit{nD-StrongSperner} problems we observed that \textit{Kleene logic}
provides a natural extension and is equivalent to its non hazard counterpart.
For other problems such as \textit{StrongSperner}, we observed that this
is not clear by current complexity assumptions and indeed proving \textbf{PPAD}-hardness
would be essential for establishing the counting hardness of \textit{PureCircuit}.
But one could ask further questions such as, could we define
complexity classes such as \textbf{HF-PPAD} and could we show that
$\textbf{HF-PPAD} = \textbf{PPAD}$.
Along similar line, we can also ask the questions as to existence of problems
\textit{HF-nD-Counting-SS} as previously stated and whether there
exists a parsimonious reduction  to its non-hazard-free variant.


Even more naturally, one can also ask further questions such as, how does hazard-freeness
affect \textbf{TFNP} as a whole. We conjecture that one can make
similar observations to the \textit{StrongTucker} problem \cite{aisenberg_2DTuckerPPA_2020},
in order to create similar classes such as \textbf{HF-PPA} as
it also uses a very similar colouring argument as the \textit{StrongSperner} argument.
Moreover we were able to observe how we can use \textit{Kleene logic} with
\textit{PLS} sub-problems like \textit{Tarski}. We believe that this shift
in perspective may also give insights to other areas such as the circuit complexity \textit{Hazard-Free Circuits}.


\subsubsection{Parsimonious Topological Representations of \textbf{PPAD} problems }


We summarise our conclusion from the section \ref{sec:sc-excess}, by arguing about
the existence of a parsimonious reduction from an \textit{nD-Counting-SS} to a
\textit{Counting-SS} problem as it would demonstrate an important stepping
stone towards its parsimonious reduction to the \textit{PureCircuit} problem.
Moreover, given our \textbf{PPAD} hierarchies, we can create the following conjecture
which combines the graph representations with the topological ones \ref{conj:conj-counting-ss-unb}


\begin{conjecturebox}{\textit{nD-CountingSS} and \textit{Unbalanced}}{conj-counting-ss-unb}
    There exist some polynomial computable function $f :\mathbb{N}^3 \to \mathbb{N}$, such that
    given $\alpha,\Delta_i, \Delta_o \in \mathbb{N}$ and $m= f(\alpha, \Delta_i, \Delta_o)$
    $$
        \textbf{PPAD}(\alpha\textit{-SUnbalanced}(\Delta_i, \Delta_o)) \subseteq
        \textbf{PPAD}(m\textit{D-CountingSS})
    $$
\end{conjecturebox}
What the above conjecture essentially describes is, given a number of known sources and the
maximum number of input/output vertices of graph, can we create some \textit{nD-CountingSS}
that is parsimonious to its graphical representation? Since we know that the problem \textit{Unbalanced}
can parsimoniously count every variant, can we also reduce \textit{Unbalanced} to a topological
problem? A reduction like that would be an important milestone for the usage topological \textit{PureCircuit} constructions
for enumerating solutions.


\subsubsection{PureCircuit Gate Extension} 

The proof for the \textbf{PPAD}-inclusion can be summarised as expressing a \textit{PureCircuit} instances
as a \textit{Brouwer} function from $[0,1]^n \to [0,1]^n$, given fixed point $x^\star$ then
we denote the following assignment: if $x^\star_i \in \{0,1\}$ then $v_i = x^{\star}_i$, otherwise
$v_i = \bot$. We argue the existence gates which exhibit any behaviour of the sort $f \in [0,1]^k \to [0,1]^m$
such that $f$ is continuous, where they give may rise to stronger combinatorial expression
than our current set.


\subsection{Software future Directions}

Given the algorithms that we proposed for our software, we argue that there are several improvements
that can be made.
For all three of the algorithmic methods we mentioned, the first and most important extension
is to treat weakly directly connected components individually. Trivially an assignment in one
weakly connected component should not affect the solution set of the other one.
We can therefore reduce the search spaces of our problems by finding
the solution sets in each graph separately. Moreover, for the
\textit{backtracking} algorithm, we can extend the above idea by considering
the Cartesian product of the solution sets, meaning:
$$
    \textbf{sol}(G) = \bigtimes_{k \in [k]}  \textbf{sol}(G_k)
$$
Where $G_i$ is the $i$th connected component.


We can also optimise the backtracking algorithm even further.
First one can observe that if our graph was acyclic, then
the permutation of the nodes with no \textit{in-degree} solutions essentially
dictate the assignments of the downstream nodes (excluding nodes incident
to purify gates). With this in mind, we can apply the following sequence
of steps to make the backtracking algorithm simpler:
\begin{enumerate}
    \item Identify the strongly connected components (SCCs) of the circuit, using polynomial
          time algorithms such as Kosaraju's Algorithm \cite{sharir_StrongconnectivityAlgorithmIts_1981} 
          or Gabow's Algorithm \cite{gabow_PathbasedDepthfirstSearch_2000}. 
    \item Treating SCCs as an individual node, identify nodes by their distances from
          base nodes (nodes with no incoming edges)
    \item Extend the backtracking algorithm so instead of the using the is start label,
          use the distance from base instead. The greater distance a node has from the
          base, the less impact it has and therefore should be postponed in later stages of the
          backtracking calls.
\end{enumerate}

We argue that the above heuristic may result in improved average times but of course
further research can be done in order to determine more suitable selection heuristics for our priority queue.

